\begin{claim}{D.32}{Proved}{The three-cusp rank test and finite pure loop groups}
Let $F=\widehat F(x,y)$, with the three named peripheral generators
$c_1=x$, $c_2=y$, and $c_3=(xy)^{-1}$.  Let $P$ be the subgroup of
$\Out(F)$ preserving each of their three procyclic conjugacy classes.
Its common peripheral multiplier is denoted by $\lambda$.
Every finite subgroup of $P$ maps injectively, by $\lambda_2$, to
$\{1,-1\}\leq\Z_2^\times$.  In particular it has order at most two,
and a nonidentity element of finite order has order two and
$\lambda_2=-1$.  These statements hold for $GT_0\leq P$ as well.

For the full three-cusp loop-outer group, allowing puncture permutations,
every finite subgroup $G$ has the explicit faithful representation
\[
 G\longrightarrow S_3\times C_2,\qquad
       g\longmapsto\bigl(\sigma_g,\lambda_2(g)\bigr).
                                                        \tag{D.32a}
\]
This is an assertion about the displayed homomorphism and abstract
finite groups.  Conjugacy of its original action to the geometric
action with this image is not asserted here. D.34 determines the
global multiplier of a surviving involution; its geometric conjugacy
remains to be proved.
\end{claim}
\begin{proof}
We first prove the closed-subgroup fact which uses all three cusps.
For a prime $p$, let $e_p\in\widehat\Z$ have $p$-coordinate one and
all other coordinates zero.  Suppose a closed subgroup $C\leq F$
has $\dim_{\F_p}H^1(C,\F_p)=2$ and contains, for each $i=1,2,3$,
a conjugate
\[
                    g_i c_i^{e_p\Z_p}g_i^{-1}.             \tag{D.32b}
\]
We prove $C=F$.

The comparison with open subgroups is the actual restriction map
\[
 \underset{U\text{ open},\ C\leq U\leq F}{\operatorname{colim}}
       H^1(U,\F_p)\ \xrightarrow{\ \sim\ } H^1(C,\F_p),
                                                        \tag{D.32c}
\]
where the transition maps are restrictions as $U$ decreases.
Indeed a continuous character of $C$ kills $C\cap V$ for some open
normal $V\triangleleft F$.  It extends to $CV$ by the quotient
$CV/V=C/(C\cap V)$.  Conversely a character on $U$ which vanishes
on $C$ vanishes after restriction to its open kernel, which contains
$C$.  This proves both surjectivity and injectivity of (D.32c).

If $C$ is open, finite-coset Schreier rewriting identifies it with
a free profinite group of rank $1+[F:C]$.  Its given first-cohomology
dimension therefore forces $[F:C]=1$.
Suppose instead that $C$ has infinite index.  The $p$-valuations
of the finite indices $[F:U]$, for open $U$ containing $C$, must
be unbounded.  To prove this, assume they were bounded and choose
$U_0$ attaining their maximum.  For $C\leq V\leq U_0$, every
further restriction between such open subgroups has prime-to-$p$
index and is injective on $H^1(-,\F_p)$: a nonzero character has
an index-$p$ kernel and cannot vanish on a subgroup of prime-to-$p$
index.  Thus (D.32c) contains the first cohomology of every such
$V$.  The indices $[F:V]$ are unbounded, since their intersection
is $C$ and $[F:C]=\infty$.  Their Schreier dimensions
$1+[F:V]$ are therefore unbounded, a contradiction.

Choose an open $U\geq C$ for which $p\mid [F:U]$.  Put
$\Delta=F(x,y)$, the named dense discrete subgroup of $F$, and
$U_0=U\cap\Delta$.  The finite covering of the pair of pants
defined by $U_0\leq\Delta$ has degree $n=[F:U]$, and its
completed fundamental group is exactly $U$: finite-index subgroups
of a discrete group carry the full induced profinite topology.
For each $g_i$ choose $h_i\in\Delta$ in its left $U$-coset and
write $g_i=u_i h_i$, with $u_i\in U$.  The subgroup
\[
 U\cap g_i\overline{\langle c_i\rangle}g_i^{-1}
  =u_i\overline{\langle h_i c_i^{d_i}h_i^{-1}\rangle}u_i^{-1}
                                                        \tag{D.32d}
\]
is consequently the completion, up to the displayed inner
conjugation, of a particular boundary subgroup in this finite
cover.  Here $d_i$ is its boundary covering degree.  Inclusion
(D.32b) proves $p\nmid d_i$.

The three selected boundaries lie over the three different original
boundaries, so they are distinct.  There is at least one more
boundary.  Otherwise each original boundary would have just one
lift, and its degree would be $n$, contrary to $p\nmid d_i$ and
$p\mid n$.  If the cover has genus $h$ and $b\geq4$ boundary
components, its presentation is
\[
 \left\langle a_1,b_1,\ldots,a_h,b_h,d_1',\ldots,d_b'
       \ \middle|\ \prod_{j=1}^h[a_j,b_j]\,
                         d_1'\cdots d_b'=1\right\rangle.
\]
Thus the only relation between its boundary homology classes over
$\F_p$ is their sum.  The three selected classes are independent,
and evaluation on them gives a surjection
$H^1(U,\F_p)\to\F_p^3$.

This remains a surjection after restriction of these characters to
$C$ and evaluation on the three elements in (D.32b).
To check the comparison explicitly, the element
$g_i c_i^{e_p}g_i^{-1}$ is the $e_p/d_i$ power of the boundary
generator in (D.32d).  Inner conjugation does not change a
character, and its evaluation is therefore multiplied by the
nonzero scalar $d_i^{-1}\in\F_p$.  We obtain three independent
characters of $C$, contradicting its dimension two.  This proves
the closed-subgroup assertion.

We now apply it to finite loop actions, retaining the precise
normalization map.  If an outer class fixes the three peripheral
classes, its exponents on $c_1,c_2,c_3$ are equal: in
$F^{\rm ab}=\widehat\Z^2$, the relation
$[c_1]+[c_2]+[c_3]=0$ gives equality of the two coordinates.
Conjugate a representative $\beta$ to have
$\beta(x)=x^\lambda$, and write
$\beta(y)=k^{-1}y^\lambda k$.  Write $[k]=(a,b)$ in the
named abelian basis.  Conjugating this representative by $x^a$
gives the normalized pair
\[
 \alpha(x)=x^\lambda,\qquad
 \alpha(y)=f^{-1}y^\lambda f,
       \qquad f=y^{-b}kx^{-a}\in F'.                     \tag{D.32e}
\]
Multiplication on the left of the conjugator by $y^{-b}$ changes
neither image of $y$ nor the automorphism.  Two normalized
representatives differ by conjugation by an element centralizing
$x^\lambda$, hence by $x^v$ by A.7.  Their two conjugators then
differ by left multiplication by a power of $y$ and right
multiplication by $x^{-v}$.  Since both are in $F'$, their two
abelian coordinates force both powers to be zero.  Normalization
is therefore unique.  A product of normalized pairs has second
coordinate $f\alpha(g)\in F'$, so (D.32e) is a homomorphic
section into $\Aut(F)$.  Its image is precisely the Belyi group
defined in \cite[p.~3120, equation (2)]{R-HS}; the argument has
not identified that entire group with its subgroup $GT_0$.

Let an element of $P$ have prime order $p$ with $\lambda_p=1$.
Its normalized automorphism has order $p$.  The proof of D.28
uses exactly the properties just established: finite order of the
normalized automorphism, the common multiplier, and preservation
of the three cusps.  Neither of the extra $GT_0$ symmetry equations
is used in its norm computation or Brown restriction argument.
In particular that proof gives a single conjugacy class of
sections in $F\rtimes C_p$ and the actual fixed subgroup
$C=C_F(C_p)$ with $\dim H^1(C,\F_p)=2$.
Each of the three peripheral sections centralizes the whole
$c_i^{e_p\Z_p}$; conjugating those sections to the specified
$C_p$ places the three groups (D.32b) in this same $C$.
The closed-subgroup assertion gives $C=F$, so the normalized
automorphism is the identity.  This contradicts its prime order.

For odd $p$, an element of order dividing $p$ in
$\Z_p^\times$ equals one, by reduction modulo $p$ and the
torsion-free group $1+p\Z_p$.  Thus no odd prime occurs in
the order of a finite subgroup of $P$.  The finite image of
$\lambda_2$ lies in $\{1,-1\}$, since $1+4\Z_2$ is
torsion-free.  Its kernel in any finite subgroup has no element
of order two by the preceding argument, and no element of odd
prime order either.  Cauchy's theorem makes that kernel trivial.
This proves the first assertion.

For a class permuting cusps, compose it with the inverse of the
corresponding positively oriented geometric permutation.  The
preceding abelian calculation gives its common multiplier too,
and composition on cusp generators proves that both $\sigma$
and $\lambda$ are homomorphisms.  On a finite subgroup, the
kernel of (D.32a) is a finite subgroup of $P$ on which
$\lambda_2=1$, hence is trivial.  All subgroups of the target
already have their geometric actions on the three-punctured
sphere: use an equilateral triple on the equator, the positive
rotational $S_3$, and reflection in that equator for the
commuting $C_2$.  This proves the claimed abstract group bound
and its attainment, while keeping the action-conjugacy question
separate through the actual map (D.32a).

The program \texttt{code/part\_c/three\_cusp\_rank\_d32.py}
enumerates the finite covers used for a validation through degree
six.  It computes the Schreier graph, every boundary cycle,
the unique boundary-homology relation, and explicit dual
characters for the selected boundaries, including the
$d_i^{-1}$ restriction scalars.  The proof above applies to
arbitrary degrees and closed subgroups; no bounded enumeration
is used in place of it.
\end{proof}
